\documentclass[../../../include/open-logic-section]{subfiles}

\begin{document}

\olfileid{sth}{cardinals}{classing}
\olsection{متناهي، \usetoken{S}{enumerable} او \usetoken{S}{nonenumerable}}

اوس چې اصلي عددونه مو وپېژندل، ښه ده د اصلي عددونو د بېلابېلو
ډولونو په اړه لږ وغږېږو؛ په ځانګړې توګه د متناهي،
!!{enumerable} او !!{nonenumerable} اصلي عددونو په اړه۔

زموږ له لومړنيو دوو پايلو څخه څرګندېږي چې متناهي اصلي عددونه به
عين متناهي ترتيبي عددونه وي؛ هماغه چې پخوا مو په
\olref[z][infinity-again]{defnomega} کښې د خپلو \emph{طبيعي
عددونو} په توګه تعريف کړل:

\begin{prop}\ollabel{finitecardisoequal}
فرض کړئ $n, m \in \omega$۔ نو $n = m$ هغه وخت او يوازې هغه وخت
چې $\cardeq{n}{m}$ وي۔
\end{prop}

\begin{proof}
له \emph{چپې خوا ښۍ خوا ته} لورے څرګند دے۔ د \emph{ښۍ خوا
چپې خوا ته} لوري د ثبوت دپاره فرض کړئ چې $\cardeq{n}{m}$، خو
$n \neq m$۔ د درې‌ګوني وېش له مخې يا $n \in m$ يا $m \in n$؛
له عموميت له کمېدو پرته $n \in m$ فرض کړئ۔ بيا $n \subsetneq m$
او يوه !!a{bijection} $f \colon m \to n$ شته، نو $m$ د
ډېډېکېنډ په معنا نامتناهي دے؛ دا له
\olref[z][infinity-again]{naturalnumbersarentinfinite} سره ټکر لري۔
\end{proof}

\begin{cor}\ollabel{naturalsarecardinals}
که $n \in \omega$ وي، نو $n$ اصلي عدد دے۔
\end{cor}

\begin{proof}
سمدستي پايله ده۔
\end{proof}
\noindent
له دې دا هم څرګندېږي چې د يوه اصلي عدد د «متناهي» يا
«نامتناهي» بللو څو معقولې معناوې سره برابرې دي:
\begin{thm}\ollabel{generalinfinitycharacter}د هر سټ $A$ دپاره لاندې
شرطونه همارزښته دي:
\begin{enumerate}
	\item\ollabel{card:notinomega} $\card{A} \notin \omega$، يعنې
	د $A$ اصلي شمېر طبيعي عدد نۀ دے؛
	\textbf{د سرچينې يادښت (OLSTH-018):} چاپي متن دلته د سټ
	پخپله طبيعي عدد نۀ کېدل وايي؛ دا د ښودل شوي شرط معنا نه ده۔
	\item\ollabel{card:omegaplus} $\omega \leq \card{A}$؛
	\item\ollabel{card:infinite} $A$ د ډېډېکېنډ په معنا نامتناهي دے۔
\end{enumerate}
\end{thm}

\begin{proof}
دا د \olref[ord-arithmetic][using-addition]{ordinfinitycharacter}،
\olref[cardinals][cardsasords]{lem:CardinalsBehaveRight} او
\olref{naturalsarecardinals} له مخې دي۔
\end{proof}

له دې سره هغه ځينې مفهومونه په \emph{رسمي ډول تعريفولے} شو چې
په \olref[sfr][][]{part} کښې مو لږ په غيررسمي توګه کارول:

\begin{defn}\ollabel{defnfinite}
وايو چې $A$ \emph{متناهي} دے که او يوازې که $\card{A}$ يو
طبيعي عدد وي، يعنې $\card{A} \in \omega$۔ کنه وايو چې $A$
\emph{نامتناهي} دے۔
\end{defn}
\noindent
خو پام وکړئ چې دا تعريف د $\ZFC$ په چوکاټ کښې وړاندې شوے دے۔
بالاخره، د هر سټ د اصلي شمېر د شتون د تضمين دپاره مو ښه ترتيب
ته اړتيا درلوده۔ په رښتيا، له ښه ترتيب پرته داسې سټ شته کېدے شي
چې نه متناهي وي او نه د ډېډېکېنډ په معنا نامتناهي۔ دې ډول مسئلې
ته به په \olref[choice][]{chap} کښې بېرته راوګرځو۔ اوس مهال پر
ښه ترتيب تکيه کوو۔

اوس له متناهي اصلي عددونو څخه نامتناهي اصلي عددونو ته راځو۔
دا دوه لومړنۍ پايلې دي:

\begin{cor}\ollabel{omegaisacardinal}
$\omega$ تر ټولو کوچنے نامتناهي اصلي عدد دے۔
\end{cor}

\begin{proof}
$\omega$ اصلي عدد دے، ځکه چې $\omega$ د ډېډېکېنډ په معنا
نامتناهي دے او که $\cardeq{\omega}{n}$ د کوم $n \in \omega$
دپاره وي، نو $n$ به هم د ډېډېکېنډ په معنا
نامتناهي وي؛ دا له
\olref[z][infinity-again]{naturalnumbersarentinfinite} سره ټکر لري۔
اوس د تعريف له مخې $\omega$ تر ټولو کوچنے نامتناهي اصلي عدد دے۔
\end{proof}

\begin{cor}
هر نامتناهي اصلي عدد يو حدي ترتيبي عدد دے۔
\end{cor}

\begin{proof}
فرض کړئ $\alpha$ يو نامتناهي تالي ترتيبي عدد وي، نو
$\alpha = \beta \ordplus 1$ د کوم $\beta$ دپاره دے۔ د
\olref{finitecardisoequal} له مخې $\beta$ هم نامتناهي دے؛ نو د
\olref[ord-arithmetic][using-addition]{ordinfinitycharacter} له مخې
$\cardeq{\beta}{\beta \ordplus 1}$۔ اوس د
\olref[cardinals][cardsasords]{lem:CardinalsBehaveRight} له مخې
$\card{\beta} = \card{\beta\ordplus 1} = \card{\alpha}$؛ نو
$\alpha \neq \card{\alpha}$۔
\end{proof}

موږ له وړاندې، په \olref[sfr][siz][enm-alt]{defn:enumerable} کښې،
دې ته پام اړولے و چې د نامتناهي سټونو تر منځ !!{enumerable} او
!!{nonenumerable} توپيرولے شو۔ هغه تعريف په طبيعي ډول لاندې
پايله راکوي:

\begin{prop}
$A$ هغه وخت او يوازې هغه وخت !!{enumerable} دے چې
$\card{A} \leq \omega$ وي؛ او $A$ هغه وخت او يوازې هغه وخت
!!{nonenumerable} دے چې $\omega < \card{A}$ وي۔
\end{prop}

\begin{proof}
د درې‌ګوني وېش له مخې دواړه ادعاوې همارزښې دي، نو بس ده دا
ثابته کړو چې $A$ هغه وخت او يوازې هغه وخت !!{enumerable} دے
چې $\card{A} \leq \omega$ وي۔ له \emph{ښۍ خوا چپې خوا ته}:
که $\card{A} \leq \omega$ وي، نو د
\olref[cardinals][cardsasords]{lem:CardinalsBehaveRight} او
\olref{omegaisacardinal} له مخې $\cardle{A}{\omega}$۔ له
\emph{چپې خوا ښۍ خوا ته}: که $A$ !!{enumerable} وي، نو د
\olref[sfr][siz][enm-alt]{defn:enumerable} له مخې درې حالتونه دي:
\begin{enumerate}
	\item که $A = \emptyset$ وي، نو د \olref{naturalsarecardinals}
	او \olref[cardinals][cardsasords]{lem:CardinalsBehaveRight} له مخې
	$\card{A} = 0 \in \omega$؛
	\item که $\cardeq{n}{A}$ وي، نو د \olref{naturalsarecardinals}
	او \olref[cardinals][cardsasords]{lem:CardinalsBehaveRight} له مخې
	$\card{A} = n \in \omega$؛
	\item که $\cardeq{\omega}{A}$ وي، نو د \olref{omegaisacardinal}
	له مخې $\card{A} = \omega$۔
\end{enumerate}
نو په ټولو حالتونو کښې $\card{A} \leq \omega$۔
\end{proof}
\noindent
په رښتيا، $\omega$ ځانګړے ځاے لري۔ که څه هم شمېر وړ ترتيبي
عددونه ډېر دي:

\begin{cor}
$\omega$ يوازينے !!{enumerable} نامتناهي اصلي عدد دے۔
\end{cor}

\begin{proof}
فرض کړئ $\cardfont{a}$ يو !!a{enumerable} نامتناهي اصلي عدد
وي۔ ځکه چې $\cardfont{a}$ نامتناهي دے، $\omega \leq
\cardfont{a}$۔ ځکه چې $\cardfont{a}$ يو !!a{enumerable} اصلي
عدد دے، $\cardfont{a} = \card{\cardfont{a}} \leq \omega$۔ نو
د درې‌ګوني وېش له مخې $\cardfont{a} = \omega$۔
\end{proof}

البته اصلي عددونه بې‌شمېره دي۔ نو پوښتنه کولے شو: \emph{څو
اصلي عددونه شته؟} لاندې پايلې ښيي چې ښايي دا پوښتنه له سره
وڅېړو۔

\begin{prop}\ollabel{unioncardinalscardinal}
که د $X$ هر غړے اصلي عدد وي، نو $\bigcup X$ هم اصلي عدد دے۔
\end{prop}

\begin{proof}
دا په اسانه کتل کېدے شي چې $\bigcup X$ ترتيبي عدد دے۔
فرض کړئ $\alpha \in \bigcup X$ يو ترتيبي عدد وي؛ نو
$\alpha \in \cardfont{b} \in X$ د کوم اصلي عدد $\cardfont{b}$
دپاره وي۔ ځکه چې $\cardfont{b}$ اصلي عدد دے،
$\cardless{\alpha}{\cardfont{b}}$۔ ځکه چې $\cardfont{b}
\subseteq \bigcup X$، نو $\cardle{\cardfont{b}}{\bigcup X}$،
او له دې $\cardneq{\alpha}{\bigcup X}$۔ په عمومي کولو سره
$\bigcup X$ اصلي عدد دے۔
\end{proof}

\begin{thm}\ollabel{lem:NoLargestCardinal}
تر ټولو لوے اصلي عدد نشته۔
\end{thm}

\begin{proof}
د هر اصلي عدد $\cardfont{a}$ دپاره د کانتور قضيه
(\olref[sfr][siz][car]{thm:cantor}) او
\olref[cardinals][cardsasords]{lem:CardinalsExist} دا راکوي چې
$\cardfont{a} < \card{\Pow{\cardfont{a}}}$۔
\end{proof}

\begin{thm}
د ټولو اصلي عددونو سټ نشته۔
\end{thm}

\begin{proof}
د تناقض دپاره فرض کړئ $C = \Setabs{\cardfont{a}}{\cardfont{a} \text{
يو اصلي عدد دے}}$۔ اوس د \olref{unioncardinalscardinal} له مخې
$\bigcup C$ اصلي عدد دے، نو د \olref{lem:NoLargestCardinal} له مخې
يو اصلي عدد $\cardfont{b} > \bigcup C$ شته۔ د تعريف له مخې
$\cardfont{b} \in C$، نو $\cardfont{b} \subseteq \bigcup{C}$،
او له دې $\cardfont{b} \leq \bigcup C$؛ دا تناقض دے۔
\end{proof}

دا د رسل له تناقض او بورالي--فورتي دواړو سره پرتله کړئ۔

\end{document}
